\documentclass[a4paper,14pt]{extarticle}
\usepackage{suai-report}

\suaisetup{
  department   = 41,
  teacher-post = Старший преподаватель,
  teacher      = В. В. Боженко,
  type         = Отчет о лабораторной работе,
  number       = 2,
  title        = Исследовательский анализ данных,
  course       = Введение в анализ данных,
  group        = 4414,
  student      = Г. В. Моисеев,
  date         = today,
}

\begin{document}

\maketitle
\suaitoc

\suaiintro

Цель работы — изучение связи между признаками двумерного набора данных, визуализация данных.

\suaitasks
загрузить набор данных, оценить его методами info и describe и выполнить предобработку
построить матрицу диаграмм рассеяния и диаграмму рассеяния по категориям
построить гистограммы числовых признаков с подобранным количеством интервалов
оценить ковариацию и корреляцию признаков и построить тепловую карту корреляции
построить графики по заданиям варианта
построить hexagonal binning plot и диаграмму размаха
добавить категорию по числовому признаку и построить по ней диаграмму размаха
построить диаграммы размаха по другим категориям

\section{Вариант задания}

Номер в списке группы — 16. В методических указаниях 15 вариантов, поэтому номер взят по кругу и выполнены задания варианта 1.

Набор данных drivers.csv содержит информацию о поездках в такси за 2016 год, каждая строка описывает одну поездку. Всего в файле 1099 поездок. Этот же набор данных использовался в лабораторной работе 1. Столбцы набора данных приведены в~\tabref{columns}.

\suaitable[columns]{Столбцы набора данных drivers.csv}
Столбец | Содержание
START\_DATE* | дата и время начала поездки
END\_DATE* | дата и время окончания поездки
CATEGORY* | категория поездки: Business (деловая) или Personal (личная)
START* | место начала поездки
STOP* | место окончания поездки
MILES* | пройденное расстояние, мили
PURPOSE* | цель поездки
time | длительность поездки, минуты
speed | средняя скорость, мили в час
price | стоимость поездки

Задания варианта:
\suainum
seaborn: по группировке по CATEGORY и количеству поездок каждой цели PURPOSE, отфильтрованной по количеству поездок больше 2, построить столбчатую диаграмму
pandas: по сводной таблице отобразить среднее количество пройденных миль по каждой цели поездки маркерами в виде звёздочки зелёного цвета размером 18
matplotlib: построить круговую диаграмму долей целей поездки, убрав цели с количеством поездок меньше 5

\section{Ход работы}

Работа выполнена в Jupyter-ноутбуке ЛР2\_Моисеев\_16.ipynb. Ссылка на ноутбук: \url{https://github.com/Echways/ItDA/tree/main/lab-2}.

\subsection{Загрузка и обзор данных}

Набор данных загружен методом read\_csv библиотеки pandas. Результат метода info приведён на~\figref{info_raw}.

\suaiimg[\textwidth]{info_raw}{Общая информация о данных до обработки}

В таблице 1099 строк и 10 столбцов. Пропуски есть только в PURPOSE*: заполнено 598 значений из 1099. Даты прочитаны как строки (object), столбцы time и price имеют тип float64.

Статистика по числовым столбцам, полученная методом describe, показана на~\figref{describe_raw}.

\suaiimg[\textwidth]{describe_raw}{Статистика по числовым столбцам до обработки}

Среднее и максимум скорости равны inf, стандартное отклонение не вычислено. Минимальная длительность time равна 0: скорость получена делением расстояния на длительность, и при нулевой длительности получается деление на ноль.

\subsection{Предобработка}

Предобработка набора данных выполнялась в лабораторной работе 1. В этой работе те же шаги повторены, чтобы ноутбук работал независимо, и добавлена проверка скорости, так как столбец speed участвует в анализе. Результаты проверок показаны на~\figref{checks}.

\suaiimg[\textwidth]{checks}{Проверка скорости и явных дубликатов}

Бесконечная скорость получена у четырёх поездок с нулевой длительностью. Ещё у 15 поездок скорость превышает 100 миль в час, что для такси нереально. Эти поездки показаны на~\figref{too_fast}.

\suaiimg[\textwidth]{too_fast}{Поездки со скоростью выше 100 миль в час}

Почти у всех таких поездок длительность составляет 1–8 минут при расстоянии 6–33 мили, например 15,1 мили за 1 минуту. Длительность записана с ошибкой, а вместе с ней ошибочны скорость и стоимость, которая, как показано в разделе~\ref{sec:correlation}, вычисляется по длительности. Восстановить эти значения нельзя, поэтому 19 поездок удалены. Бесконечная скорость также больше 100, поэтому оба случая отсекаются одним условием. Явных дубликатов нет, а неявных дубликатов в CATEGORY и PURPOSE, как показала лабораторная работа 1, тоже нет.

Шаги предобработки оформлены отдельными функциями, каждая из которых принимает датафрейм и возвращает новый. Функции объединены в цепочку методом pipe (\lstref{preprocessing}):
\suailist
strip\_column\_marks удаляет символ «*» в конце названий столбцов
drop\_impossible\_trips удаляет поездки со скоростью выше 100 миль в час
fill\_missing\_purpose заменяет 501 пропуск в цели поездки значением «не определена», так как удаление строк привело бы к потере почти половины данных
convert\_types переводит даты в datetime, а time и price — в int64

\begin{code}[Python, preprocessing]{Предобработка данных}
MAX_SPEED = 100
UNKNOWN_PURPOSE = 'не определена'


def strip_column_marks(trips: pd.DataFrame) -> pd.DataFrame:
    return trips.rename(columns=lambda name: name.rstrip('*'))


def drop_impossible_trips(trips: pd.DataFrame) -> pd.DataFrame:
    return trips[trips['speed'] <= MAX_SPEED].reset_index(drop=True)


def fill_missing_purpose(trips: pd.DataFrame) -> pd.DataFrame:
    return trips.fillna({'PURPOSE': UNKNOWN_PURPOSE})


def convert_types(trips: pd.DataFrame) -> pd.DataFrame:
    dates = {column: pd.to_datetime(trips[column])
             for column in ['START_DATE', 'END_DATE']}
    return trips.assign(**dates).astype({'time': 'int64', 'price': 'int64'})


trips = (raw.pipe(strip_column_marks)
            .pipe(drop_impossible_trips)
            .pipe(fill_missing_purpose)
            .pipe(convert_types))
\end{code}

Результат предобработки показан на~\figref{info_clean}, статистика после обработки — на~\figref{describe_clean}.

\suaiimg[0.6\textwidth]{info_clean}{Общая информация о данных после обработки}

\suaiimg[\textwidth]{describe_clean}{Статистика по числовым столбцам после обработки}

После обработки осталось 1080 поездок, пропусков нет. Расстояние MILES лежит в пределах от 0,5 до 310,3 мили, медиана 6, среднее 10,6: среднее почти вдвое больше медианы из-за нескольких дальних поездок. Длительность time — от 1 до 336 минут, медиана 17. Скорость speed — от 3,9 до 91,5 мили в час, медиана 21,1, стандартное отклонение 12,5. Стоимость price — от 412 до 25~569, медиана 1611.

\subsection{Диаграммы рассеяния}

Матрица диаграмм рассеяния для числовых признаков построена функцией pairplot библиотеки seaborn, точки раскрашены по категории поездки (\figref{pairplot}).

\suaiimg[0.9\textwidth]{pairplot}{Матрица диаграмм рассеяния по категориям поездки}

Точки time и price лежат на одной прямой: стоимость полностью определяется длительностью. Расстояние MILES связано с time и price прямой зависимостью — чем дальше поездка, тем она дольше и дороже, при этом разброс растёт с расстоянием. Скорость speed слабо связана с длительностью и стоимостью. У поездок до 5 миль скорость разбросана от 6 до 68 миль в час, а у поездок длиннее 50 миль она не опускается ниже 25 миль в час, так как это поездки по трассе. Личные поездки лежат в той же области, что и деловые, и отдельной группы не образуют.

Большинство точек сжато в левом нижнем углу, поэтому диаграмма рассеяния расстояния и длительности по целям поездки построена в логарифмическом масштабе (\figref{scatter_purpose}). Поездки без цели на ней не показаны: их почти половина, и они закрывают остальные точки.

\suaiimg[\textwidth]{scatter_purpose}{Расстояние и длительность поездок по целям}

Точки вытянуты вдоль диагонали. Цели занимают разные участки: Errand/Supplies и Meal/Entertain в основном расположены до 10 миль, а дальше 30 миль находятся почти только Customer Visit и Meeting.

\subsection{Гистограммы}

Количество интервалов подобрано по правилу Фридмана — Диакониса, ширина интервала вычисляется по~\formref{fd}:
\suaieq[fd]{h = \frac{2 \cdot IQR}{\sqrt[3]{n}}}
h | ширина интервала
IQR | межквартильный размах
n | количество значений

Правило опирается на межквартильный размах, поэтому выбросы не влияют на ширину интервала. Гистограммы построены в обычной и в логарифмической шкале (\lstref{histograms}). Для логарифмической шкалы правило применено к логарифмам значений.

\begin{code}[Python, histograms]{Построение гистограмм}
for column, linear_ax, log_ax in zip(NUMERIC, *axes):
    values = trips[column]
    linear_bins = np.histogram_bin_edges(values, bins='fd')
    log_bins = 10 ** np.histogram_bin_edges(np.log10(values), bins='fd')

    linear_ax.hist(values, bins=linear_bins, linewidth=0)
    log_ax.hist(values, bins=log_bins, linewidth=0)
    log_ax.set_xscale('log')
\end{code}

Полученные гистограммы показаны на~\figref{histograms}.

\suaiimg[\textwidth]{histograms}{Гистограммы числовых признаков в обычной (сверху) и логарифмической (снизу) шкале}

В обычной шкале для MILES, time и price правило даёт 96–207 интервалов, но почти все поездки попадают в первые из них, а справа тянется длинный редкий хвост: 25 поездок длиннее 50 миль при медиане 6 миль. Такие гистограммы показывают асимметрию, но не форму основной массы значений. У speed хвост короче, и 34 интервалов хватает: распределение одновершинное с пиком около 20 миль в час и скошено вправо.

В логарифмической шкале всем признакам достаточно 25–29 интервалов, распределения принимают форму колокола и близки к логнормальному. У time слева видны пустые интервалы: длительность записана целыми минутами. Поэтому для этих данных выбрана логарифмическая шкала с 25–29 интервалами.

\subsection{Корреляция и ковариация}
\label{sec:correlation}

Ковариация показывает, изменяются ли две величины в одну сторону, но её значение зависит от единиц измерения. Коэффициент корреляции Пирсона нормирует ковариацию и вычисляется по~\formref{pearson}:
\suaieq[pearson]{r_{xy} = \frac{\operatorname{cov}(x, y)}{\sigma_x \sigma_y}}
\operatorname{cov}(x, y) | ковариация признаков x и y
\sigma_x, \sigma_y | стандартные отклонения признаков

Коэффициент Пирсона оценивает линейную связь. Коэффициент Спирмена вычисляется так же, но по рангам значений, поэтому показывает монотонную связь любой формы и меньше чувствителен к выбросам. Матрица ковариаций приведена на~\figref{covariance}.

\suaiimg[\textwidth]{covariance}{Матрица ковариаций}

Все ковариации положительные: ни один признак не убывает при росте другого. Величина ковариации определяется масштабом признаков, у пар со стоимостью она измеряется десятками тысяч, поэтому о силе связи по этой таблице судить нельзя. Тепловые карты корреляции по Пирсону и Спирмену показаны на~\figref{heatmap}.

\suaiimg[\textwidth]{heatmap}{Тепловые карты корреляции}

Корреляция time и price равна 1,00: связь линейная и полная. MILES сильно связан с time и price, коэффициенты 0,85–0,87 по обоим методам. Связь MILES и speed средняя: 0,48 по Пирсону и 0,62 по Спирмену. Коэффициент Спирмена выше, значит связь скорее монотонная, чем линейная. Связь speed с time и price слабая, 0,18–0,24: долгая поездка может быть как дальней, так и медленной. Отрицательных корреляций нет.

Так как корреляция time и price равна единице, по методу наименьших квадратов найдена формула стоимости (\figref{price_formula}).

\suaiimg[\textwidth]{price_formula}{Линейная зависимость стоимости от длительности}

Стоимость складывается из примерно 335 за посадку и 75 за минуту поездки, отклонение от формулы не превышает 29. От расстояния стоимость зависит только через длительность, поэтому price не несёт новой информации по сравнению с time.

\subsection{Задания варианта}

Код задания 1 приведён в~\lstref{task1}, результат — на~\figref{task1}.

\begin{code}[Python, task1]{Задание 1}
purpose_counts = (trips.groupby(['CATEGORY', 'PURPOSE'])
                       .size()
                       .reset_index(name='count')
                       .query('count > 2'))

fig, ax = plt.subplots(figsize=(14, 6))
sns.barplot(data=purpose_counts, x='PURPOSE', y='count',
            hue='CATEGORY', ax=ax)
\end{code}

\suaiimg[\textwidth]{task1}{Количество поездок по целям и категориям}

Фильтр убрал Airport/Travel (2 поездки), Charity (\$) и Commute (по одной поездке). Почти все поездки деловые: 1003 из 1080. У деловых поездок чаще всего указаны цели Meeting (178), Meal/Entertain (145) и Errand/Supplies (110), но больше всего поездок без цели (415). У личных поездок цель почти никогда не указана: 71 из 77, из указанных целей после фильтра остался только Moving (4).

Код задания 2 приведён в~\lstref{task2}, результат — на~\figref{task2}.

\begin{code}[Python, task2]{Задание 2}
mean_miles = trips.pivot_table(index='PURPOSE', values='MILES',
                               aggfunc='mean')

ax = mean_miles.plot(style='*', color='green', markersize=18,
                     figsize=(14, 7), xticks=range(len(mean_miles)),
                     rot=90)
\end{code}

\suaiimg[\textwidth]{task2}{Среднее число пройденных миль по целям поездки}

Наибольшее среднее расстояние у Commute — 180,2 мили, но оно посчитано по одной поездке. Остальные средние лежат в пределах 4–22 миль. Дальше всего в среднем ездят к клиентам (Customer Visit, 21,9 мили) и на встречи (Meeting, 15,6), ближе всего — по поручениям (Errand/Supplies, 4,1), при переезде (Moving, 4,6) и на обед (Meal/Entertain, 5,7). Среднее чувствительно к дальним поездкам: медиана Customer Visit составляет всего 8,3 мили.

Код задания 3 приведён в~\lstref{task3}, результат — на~\figref{task3}.

\begin{code}[Python, task3]{Задание 3}
purpose_share = (trips['PURPOSE'].value_counts()
                                 .loc[lambda counts: counts >= MIN_TRIPS]
                                 .sort_index())

plt.figure(figsize=(20, 10))
plt.pie(purpose_share, labels=purpose_share.index, autopct='%1.1f%%',
        textprops={'size': 'x-large'}, startangle=90)
plt.legend(fontsize=16, bbox_to_anchor=(1, 1))
\end{code}

\suaiimg[0.9\textwidth]{task3}{Распределение целей поездки}

Из диаграммы убраны Moving, Airport/Travel, Charity (\$) и Commute — вместе 8 поездок. Почти половину диаграммы (45,3~\%) занимают поездки без цели. Среди известных целей чаще всего встречаются Meeting (16,6~\%), Meal/Entertain (13,5~\%) и Errand/Supplies (10,3~\%). На Customer Visit приходится 8,6~\%, на Temporary Site — 4,1~\%, на Between Offices — 1,6~\%.

\subsection{Hexagonal binning plot}

Плотность поездок по расстоянию и длительности показана на~\figref{hexbin}. Чтобы редкие дальние поездки не растягивали оси, выбраны поездки до 30 миль и до 60 минут — 1016 из 1080 (94~\%). Пустые шестиугольники не закрашены.

\suaiimg[0.75\textwidth]{hexbin}{Плотность поездок по расстоянию и длительности}

Самые плотные ячейки соответствуют поездкам на 1–4 мили длительностью 5–10 минут — это типичная поездка по городу. Далее плотность вытягивается полосой вдоль диагонали: на 10 миль половина поездок тратит от 18 до 29 минут. С расстоянием полоса расширяется, то есть время поездки на одно и то же расстояние разбросано сильнее.

\subsection{Диаграмма размаха стоимости}

Диаграмма размаха стоимости построена средствами matplotlib (\figref{boxplot_price}), квартили и граница выбросов приведены на~\figref{price_quartiles}.

\suaiimg[\textwidth]{boxplot_price}{Диаграмма размаха стоимости поездки}

\suaiimg[\textwidth]{price_quartiles}{Квартили и граница выбросов стоимости}

Половина поездок стоит от 1087 до 2437, медиана 1611 расположена ближе к левому краю ящика. Левый ус короткий и заканчивается на минимуме 412, выбросов снизу нет. Справа за границей 4462 лежат 59 выбросов (5,5~\% поездок), самые дорогие — две поездки около 25~000. Это не ошибки: 53 из 59 выбросов — поездки дальше 15 миль, а две самые дорогие поездки длились около 5,5 часа. Распределение стоимости сильно скошено вправо.

\subsection{Категория дальности поездки}

По столбцу MILES добавлена категория дальности: короткая — до 5 миль, средняя — от 5 до 15 миль, дальняя — больше 15 миль (\lstref{distance}). Границы близки к медиане (6 миль) и 90-му процентилю (17 миль) расстояния.

\begin{code}[Python, distance]{Категория дальности поездки}
DISTANCE_BINS = [0, 5, 15, np.inf]
DISTANCE_NAMES = ['короткая', 'средняя', 'дальняя']

distance = pd.cut(trips['MILES'], bins=DISTANCE_BINS,
                  labels=DISTANCE_NAMES)
trips = trips.assign(distance=distance)
\end{code}

Количество поездок в каждой категории показано на~\figref{distance_counts}. Коротких и средних поездок почти поровну (467 и 469), дальних — 144.

\suaiimg[\textwidth]{distance_counts}{Количество поездок по категориям дальности}

Диаграммы размаха построены средствами seaborn для расстояния и, для сравнения, для стоимости (\figref{boxplot_distance}). Обе оси значений логарифмические.

\suaiimg[\textwidth]{boxplot_distance}{Расстояние и стоимость по категориям дальности}

Ящики расстояния не пересекаются, так заданы категории. Медианы составляют 2,5, 8,6 и 21 милю. Самый большой разброс у дальних поездок: от 15 до 310 миль, выше 70 миль расположены выбросы. Медианы стоимости растут так же: 1087, 1913 и 3864. Ящики стоимости не пересекаются, но усы и выбросы соседних категорий перекрываются: 18 коротких поездок дороже медианы средних. Стоимость вычисляется по времени, а короткая поездка может затянуться, например, из-за пробок.

\subsection{Диаграммы размаха по другим категориям}

Диаграмма размаха скорости по категориям поездки построена средствами pandas (\figref{boxplot_category}).

\suaiimg[0.7\textwidth]{boxplot_category}{Скорость по категориям поездки}

Медиана скорости деловых поездок 21,3 мили в час, личных — 18,9. Половина деловых поездок идёт со скоростью 15–29 миль в час, личных — 14–24. У деловых поездок много выбросов вверх, до 91,5 мили в час, — это дальние поездки по трассе. Самая высокая скорость личной поездки — 58,4 мили в час. Разница между категориями небольшая, к тому же личных поездок всего 77.

Диаграмма размаха расстояния по целям поездки построена средствами seaborn (\figref{boxplot_purpose}). Показаны цели, у которых не меньше 5 поездок, в порядке возрастания медианы расстояния.

\suaiimg[\textwidth]{boxplot_purpose}{Расстояние по целям поездки}

Цели делятся на две группы. Короткие поездки — Errand/Supplies (медиана 3,1 мили) и Meal/Entertain (4,1). Длинные — Temporary Site (8,0), Customer Visit (8,3), Meeting (10,2) и Between Offices (11,8). У Customer Visit и Meeting много выбросов справа, до 100–300 миль: к клиентам и на встречи иногда ездят в другой город. Поэтому в задании 2 средние у этих целей получились намного больше медиан. У поездок без цели (медиана 4,9) самый широкий ящик, так как в эту группу попали поездки с разными целями.

\suaiconclusion

В ходе работы выполнен исследовательский анализ набора данных drivers.csv, содержащего 1099 поездок в такси за 2016 год: расстояние, длительность, скорость, стоимость, категорию и цель каждой поездки.

При предобработке удалены 19 поездок с ошибочной длительностью: у 4 она нулевая, у 15 скорость превышает 100 миль в час. Из названий столбцов удалён символ «*», пропуски в цели поездки заменены значением «не определена», даты переведены в datetime, длительность и стоимость — в int64. Осталось 1080 поездок.

Расстояние, длительность и стоимость распределены с длинным правым хвостом и близки к логнормальному распределению, поэтому гистограммы и часть диаграмм построены в логарифмической шкале. Медианная поездка — 6 миль и 17 минут, чаще всего встречаются поездки на 1–4 мили длительностью 5–10 минут.

Корреляционный анализ показал, что стоимость линейно зависит от длительности: около 335 за посадку и 75 за минуту, коэффициент корреляции 1,00. Расстояние сильно связано с длительностью и стоимостью (0,85–0,87) и умеренно — со скоростью (0,48 по Пирсону, 0,62 по Спирмену): дальние поездки проходят по трассе и быстрее. Отрицательных связей нет.

По заданиям варианта: 93~\% поездок деловые, у 45~\% поездок цель не указана, среди известных целей чаще всего встречаются встречи, обеды и поручения. Диаграммы размаха показали, что цели делятся на короткие (поручения и обеды, медиана 3–4 мили) и длинные (встречи, визиты к клиентам, поездки между офисами, медиана 8–12 миль). Средние расстояния у встреч и визитов к клиентам завышены редкими междугородними поездками, медиана описывает эти цели точнее. Скорость деловых и личных поездок различается слабо. Выбросы стоимости в основном соответствуют дальним поездкам и ошибками не являются.

\suaisources
Учебно-методические указания к выполнению лабораторной работы №2 по дисциплине «Введение в анализ данных».

\end{document}
